% Part: methods
% Chapter: proofs
% Section: reading-proofs

\documentclass[../../../include/open-logic-section]{subfiles}

\begin{document}

\olfileid{mth}{prf}{rea}
\olsection{প্রমাণ পড়া}

পাঠ্যবই বা প্রবন্ধের প্রমাণে আমাদের উদাহরণগুলির মতো
সব খুঁটিনাটি সচরাচর লেখা থাকে না। লেখক কখন
ক্ষেত্রবিচার করছেন, কখন পরোক্ষ প্রমাণ দিচ্ছেন,
অথবা কোনো সংজ্ঞা ব্যবহার করছেন—অনেক সময় তা
আলাদা করে বলেন না। তাই পাঠ্যবইয়ের প্রমাণ বুঝতে
সেই ধাপগুলি নিজেকেই পূরণ করতে হয়। এতে প্রমাণ
লেখার নানা কৌশল চেনার অনুশীলনও হয়। একটি উদাহরণ দেখি।

\begin{prop}[শোষণ]
সব সেট $A$, $B$-র জন্য,
\[
A \cap (A \cup B) = A
\]
\end{prop}

\begin{proof}
$z \in A \cap (A \cup B)$ হলে $z \in A$;
তাই $A \cap (A \cup B) \subseteq A$। এবার
$z \in A$ ধরি। তাহলে $z \in A \cup B$-ও,
আর তাই $z \in A \cap (A \cup B)$-ও।
\end{proof}

শোষণ-সূত্রের উপরের প্রমাণটি খুব সংক্ষিপ্ত। কোন
সংজ্ঞা ব্যবহার হয়েছে বা আগে ``দেখাতে হবে''
কী, তা বলা নেই। এবার বিস্তারিত খুলে দেখি।
প্রমাণিত প্রস্তাবটি যেকোনো সেট $A$, $B$ নিয়ে
সার্বিক দাবি; প্রমাণে $A$ বা $B$ হলো যেকোনো
সেটের চলক। যে সার্বিক দাবিগুলি প্রতিষ্ঠিত হয়,
সেগুলিই সেটের অভেদ প্রমাণে দরকার: সমতার বাম
দিকের প্রত্যেক !!{element} ডান দিকেরও
!!a{element}, এবং উল্টোটিও।

\begin{quote}
``$z \in A \cap (A \cup B)$ হলে $z \in A$;
তাই $A \cap (A \cup B) \subseteq A$।''
\end{quote}

এটি অভেদের প্রমাণের প্রথমার্ধ: বাম দিকের
!!a{element} যেকোনো~$z$ ডান দিকেরও
!!a{element}, অর্থাৎ
$A \cap (A \cup B) \subseteq A$। ধরো,
$z \in A \cap (A \cup B)$। $z$ দুটি সেটের ছেদের
!!{element} হওয়া মানে উভয়েরই !!{element}
হওয়া। তাই $z \in A$ এবং $z \in A \cup B$।
বিশেষত $z \in A$—এটিই দেখাতে চেয়েছিলাম।
প্রথমার্ধে আর কিছু বাকি নেই; সুতরাং বাকি প্রমাণে
দ্বিতীয়ার্ধ, অর্থাৎ $A \subseteq A \cap
(A \cup B)$, দেখানো হবে।

\begin{quote}
``এবার $z \in A$ ধরি। তাহলে $z \in A \cup B$-ও,
আর তাই $z \in A \cap (A \cup B)$-ও।''
\end{quote}

প্রথমে $z \in A$ ধরি, কারণ প্রত্যেক~$z$-এর
জন্য $z \in A$ হলে $z \in A \cap (A \cup B)$
দেখাচ্ছি। $z \in A \cap (A \cup B)$ দেখাতে
``$\cap$''-র সংজ্ঞা অনুযায়ী (i) $z \in A$
এবং (ii) $z \in A \cup B$ দেখাতে হবে। (i)
আমাদের পূর্বধারণাই, তাই আবার প্রমাণ বা উল্লেখ
করতে হয়নি। (ii)-এর জন্য মনে করো, $z$ সংযোগের
!!a{element} তখনই, যখন সেটদুটির অন্তত একটির
!!{element}। $z \in A$, আর $A \cup B$ হলো
$A$ ও $B$-র সংযোগ; তাই এখানে শর্তটি পূরণ
হয়। ফলে $z \in A \cup B$। (i) ও (ii)
দুটিই দেখিয়েছি; অতএব ``$\cap$''-র সংজ্ঞায়
$z \in A \cap (A \cup B)$। সংক্ষিপ্ত প্রমাণে
সংজ্ঞাগুলির কথা লেখা নেই; ধরে নেওয়া হয়েছে
পাঠকের জানা। জানা না থাকলে আবার দেখে নিতে
হবে। তারপর চিনতে হবে, $z \in A$ থেকে কেন
$z \in A \cup B$, আর $z \in A$ এবং
$z \in A \cup B$ থেকে কেন
$z \in A \cap (A \cup B)$ আসে।

সব ধাপ স্পষ্ট করে উপরের প্রমাণের আরেকটি রূপ:
\begin{proof}{}
[সেটের $=$-র সংজ্ঞায় $A \cap (A \cup B) = A$
  দেখাতে (a) $A \cap (A \cup B) \subseteq A$
  % BN-SRC-481: the second inclusion must run from A into the intersection.
  এবং (b) $A \subseteq A \cap (A \cup B)$
  দেখাতে হবে। (a): $\subseteq$-র সংজ্ঞায়
  $z \in A \cap (A \cup B)$ হলে $z \in A$
  দেখাতে হবে।] $z \in A \cap (A \cup B)$ হলে
  $z \in A$ [$\cap$-র সংজ্ঞায়
  $z \in A \cap (A \cup B)$ তখনই, যখন
  $z \in A$ এবং $z \in A \cup B$]; তাই
  $A \cap (A \cup B) \subseteq A$।
  [(b): $\subseteq$-র সংজ্ঞায় $z \in A$ হলে
  $z \in A \cap (A \cup B)$ দেখাতে হবে।]
  এবার [(1)] $z \in A$ ধরি। তাহলে
  [(2)] $z \in A \cup B$-ও [$z \in A$ অর্থাৎ
  (1) থেকে $z \in A$ অথবা $z \in B$;
  $\cup$-র সংজ্ঞায় এর অর্থ $z \in A \cup B$]।
  অতএব $z \in A \cap (A \cup B)$
  [$\cap$-র সংজ্ঞায় $z \in A$, অর্থাৎ (1),
  % BN-SRC-482: remove the unmatched closing parenthesis in the source.
  এবং $z \in A \cup B$, অর্থাৎ (2), দরকার]।
\end{proof}

\begin{prob}
নিচের $A \cup (A \cap B) = A$ প্রমাণটি বিস্তারিত
করো। কোন কোন অনুমানের রীতি ব্যবহৃত হয়েছে,
প্রতিটি ধাপ কোন পূর্বধারণা বা আগের প্রতিষ্ঠিত
বক্তব্য থেকে এসেছে, এবং কোথায় কোন সংজ্ঞা
লাগছে—সব উল্লেখ করো।
\begin{proof}
  $z \in A \cup (A \cap B)$ হলে $z \in A$
  অথবা $z \in A \cap B$। $z \in A \cap B$
  হলে $z \in A$। যেকোনো $z \in A$-র জন্য
  $\in A \cup (A \cap B)$-ও।
\end{proof}
\end{prob}

\end{document}
